\documentclass[10pt,a4paper]{article}
\usepackage{nexusS5}
\setnexusheader{Entrée en Première générale — Spécialité mathématiques — Stage de pré-rentrée}{Séance 5/5 — Malek KHADHRANI}
\begin{document}
\nexuscover{SÉANCE 5 --- RÉINVESTIR, CONSOLIDER, PRÉPARER LA RENTRÉE}{Livret de travail}{Malek KHADHRANI}{Substitution rigoureuse et vecteurs}{Entrée en Première générale — Spécialité mathématiques --- Mathématiques --- 1\,h\,15 de travail, puis 45 minutes d'évaluation dans un document séparé}
\begin{goalbox}
\textbf{Objectif personnel de cette séance.} Ralentir lorsque la question paraît familière et contrôler la règle activée — conseil du dossier — en sécurisant la substitution d'une valeur négative puis le test de colinéarité.
\par\medskip
\textbf{Ce livret couvre uniquement les 75 premières minutes.} L'évaluation finale est remise ensuite, dans un document distinct.
\end{goalbox}
\begin{methodbox}
\textbf{Le fil des cinq séances}
\begin{itemize}[leftmargin=13mm,labelsep=3mm,itemsep=1.5pt]
\item[\textbf{\color{Navy}S1}] Calcul algébrique, inéquations et transition vers le second degré ;
\item[\textbf{\color{Navy}S2}] Fonctions, fonctions de référence et première approche de la dérivation ;
\item[\textbf{\color{Navy}S3}] Vecteurs, droites et transition vers le produit scalaire ;
\item[\textbf{\color{Navy}S4}] Pourcentages, événements et probabilités conditionnelles ;
\item[\textbf{\color{Navy}S5}] aujourd'hui : réactiver, consolider deux axes personnels, faire le pont vers la Première spécialité, puis mesurer.
\end{itemize}
\end{methodbox}
\sectiontitle{Feuille de route des 75 minutes}
\rowcolors{2}{SoftBlue}{white}
\par\noindent\begin{tabularx}{\linewidth}{>{\bfseries}p{20mm} X X}
\rowcolor{Navy}\color{white}Temps & \color{white}Travail & \color{white}Preuve attendue \\
0--10 min & Réactiver les cinq blocs du stage (algèbre, fonctions, vecteurs, droites, pourcentages) avant la phase de consolidation. & Cinq réponses écrites et le contrôle utilisé. \\
10--35 min & Calculer une image, en particulier pour une valeur négative & Trois réussites autonomes sur l'axe prioritaire. \\
35--55 min & Colinéarité de deux vecteurs et alignement de trois points & Deux réussites autonomes sur le second axe. \\
55--70 min & De la Seconde à la Première : la suite et le taux de variation & Une production reliant un acquis de Seconde à un attendu de Première spécialité. \\
70--75 min & Cadrage méthodologique, sans aucune information sur le contenu de l'évaluation. & Deux engagements écrits avant l'évaluation. \\
\end{tabularx}
\begin{graybox}\small\textbf{Règle de la séance.} J'écris ma réponse \emph{et} le contrôle qui la rend défendable. Une réponse sans contrôle reste une hypothèse. La ligne « Certitude » se remplit \emph{avant} toute vérification.\end{graybox}
\par\vspace{2mm}
\phasehead{Phase 1 --- Réactivation}{0--10 min}{Réactiver les cinq blocs du stage (algèbre, fonctions, vecteurs, droites, pourcentages) avant la phase de consolidation.}
Sept minutes seules, sans document, puis trois minutes de mise en commun. Écrire la réponse \emph{et} le contrôle utilisé.
\begin{enumerate}
\item Factoriser $9x^2 - 25$.
\shortlines{2}
\item Résoudre $-4x + 10 \ge 0$ et donner l'ensemble des solutions sous forme d'intervalle.
\shortlines{2}
\item Pour $f(x)=x^2-5x$, calculer $f(-3)$ en écrivant la substitution complète.
\shortlines{2}
\item $A(-2\,;5)$ et $B(3\,;-1)$ : donner les coordonnées de $\vec{AB}$, puis le coefficient directeur de $(AB)$.
\shortlines{2}
\item Un prix passe de $250$ TND à $210$ TND. Calculer le taux d'évolution en pourcentage.
\shortlines{2}
\end{enumerate}
\certline
\aideline
\needspace{62mm}\par\vspace{3mm}
\phasehead{Phase 2 --- Consolidation, axe prioritaire}{10--35 min}{Calculer une image, en particulier pour une valeur négative}
\begin{methodbox}
\textbf{Substituer, c'est remplacer la lettre par le nombre \emph{entre parenthèses}.}
$(-3)^2 = 9$ alors que $-3^2 = -9$ : dans la seconde écriture, le carré ne porte que sur $3$.
\textbf{Contrôle :} réécrire l'expression complète avec les parenthèses avant d'effectuer le moindre calcul.
\end{methodbox}
\begin{graybox}
\textbf{Exemple guidé.} Pour $f(x) = x^2 - 4x$, calculer $f(-3)$.
\begin{enumerate}
\item Écrire la substitution complète : $f(-3) = (-3)^2 - 4\times(-3)$.
\item Calculer chaque terme : $9$ et $+12$.
\item Conclure : $f(-3) = 21$.
\item Contrôle : les deux termes sont positifs, le résultat ne peut pas être négatif.
\end{enumerate}
\end{graybox}
\subsectiontitle{À toi}
\begin{enumerate}
\needspace{18mm}
\item Soit $f(x) = x^2 + 2x$. Calculer $f(-4)$, $f(0)$ et $f(3)$.
\worklines{4}
\needspace{18mm}
\item Soit $g(x) = -2x^2 + x$. Calculer $g(-3)$.
\worklines{3}
\needspace{18mm}
\item Soit $h(x) = (x - 1)^2$. Calculer $h(-4)$.
\worklines{3}
\needspace{18mm}
\item Calculer séparément $(-5)^2$ et $-5^2$, puis expliquer la différence en une phrase.
\worklines{3}
\needspace{18mm}
\item Soit $k(x) = x^3$. Calculer $k(-2)$ et justifier le signe obtenu.
\worklines{3}
\needspace{18mm}
\item Pour $f(x) = x^2 - 3x$, un élève écrit $f(-2) = -2^2 - 3\times(-2) = -4 + 6 = 2$. Nommer l'erreur, corriger et donner le contrôle qui l'aurait détectée.
\worklines{4}
\end{enumerate}
\begin{successbox}\textbf{Critère de réussite de cette phase :} trois substitutions négatives consécutives sans erreur de signe, la substitution complète étant écrite avant tout calcul.\end{successbox}
\certline
\needspace{62mm}\par\vspace{3mm}
\phasehead{Phase 3 --- Consolidation, second axe}{35--55 min}{Colinéarité de deux vecteurs et alignement de trois points}
\begin{methodbox}
\textbf{Un seul calcul.} Les vecteurs $\vec{u}(x\,;y)$ et $\vec{v}(x'\,;y')$ sont colinéaires si et seulement si
le \textbf{déterminant} $xy' - yx'$ est nul.
\par\smallskip
Trois points $A$, $B$, $C$ sont alignés si et seulement si $\vec{AB}$ et $\vec{AC}$ sont colinéaires.
\textbf{Contrôle :} deux vecteurs colinéaires ont le même coefficient directeur, lorsqu'il existe.
\end{methodbox}
\subsectiontitle{À toi}
\begin{enumerate}
\needspace{18mm}
\item Les vecteurs $\vec{u}(5\,;-2)$ et $\vec{v}(-15\,;6)$ sont-ils colinéaires ? Justifier par le calcul du déterminant.
\worklines{3}
\needspace{18mm}
\item Les vecteurs $\vec{u}(4\,;7)$ et $\vec{v}(2\,;3)$ sont-ils colinéaires ?
\worklines{3}
\needspace{18mm}
\item On donne $A(1\,;2)$, $B(4\,;8)$ et $C(-1\,;-2)$. Montrer que ces trois points sont alignés.
\worklines{4}
\needspace{18mm}
\item Déterminer $k$ pour que $\vec{u}(3\,;k)$ et $\vec{v}(9\,;15)$ soient colinéaires.
\worklines{4}
\end{enumerate}
\begin{successbox}\textbf{Critère de réussite de cette phase :} quatre tests de colinéarité corrects consécutifs, dont un alignement de trois points, le déterminant étant écrit à chaque fois.\end{successbox}
\certline
\needspace{62mm}\par\vspace{3mm}
\phasehead{Phase 4 --- Réinvestissement}{55--70 min}{Découverte guidée : construire, à partir des évolutions en pourcentage et du coefficient directeur déjà travaillés, deux notions nouvelles de Première — la suite géométrique définie par récurrence et le taux de variation moyen.}

\begin{methodbox}
\textbf{Ce qui change en Première.} Deux objets \textbf{nouveaux} s'appuient sur des acquis de Seconde.
\par\smallskip
Une \textbf{suite} est une liste ordonnée de nombres indexée par les entiers. Toutes les suites ne décrivent
pas une évolution répétée ; en revanche, \emph{certaines} d'entre elles, notamment les \textbf{suites
géométriques}, modélisent exactement cela — c'est le cas étudié ci-dessous.
\par\smallskip
Le \textbf{taux de variation} d'une fonction entre deux points de sa courbe se calcule comme le coefficient
directeur de la droite qui les joint. C'est cette lecture qui conduira, en Première, au nombre dérivé.
\end{methodbox}
\begin{warningbox}\small
\textbf{Cette phase introduit des notions de l'année à venir.} Ce que tu produis ici ne sera donc
\emph{pas} comparé à ton positionnement de début de stage : il ne s'agit pas de mesurer un progrès,
mais de préparer les premières séances de septembre.
\end{warningbox}


\subsectiontitle{A. De l'évolution à la suite — 8 min}
Un capital de $2\,000$ TND est placé à $3\,\%$ par an. On note $u_n$ le capital au bout de $n$ années.
\begin{enumerate}
\item Donner le coefficient multiplicateur annuel. \rule{35mm}{0.32pt} \quad Écrire $u_0$. \rule{28mm}{0.32pt}
\item Écrire la relation qui donne $u_{n+1}$ en fonction de $u_n$. \rule{55mm}{0.32pt}
\item Calculer $u_1$ et $u_2$ (arrondir au centime).
\shortlines{2}
\item Compléter le programme Python ci-dessous, qui affiche $u_0$ puis les cinq termes suivants.
\begin{lstlisting}
u = 2000
print(0, u)
for n in range(5):
    u = _______________________
    print(________, u)
\end{lstlisting}
\par\vspace{0.5mm}
Sans la ligne \code{print(0, u)}, quels termes le programme afficherait-il ? \rule{45mm}{0.32pt}
\item Le capital dépasse-t-il $2\,200$ TND au bout de trois ans ? Répondre par un calcul.
\shortlines{2}
\end{enumerate}

\subsectiontitle{B. Du coefficient directeur au taux de variation — 7 min}
Soit $f$ la fonction définie par $f(x) = x^2$.
\begin{enumerate}
\item Calculer le coefficient directeur de la droite passant par les points de la courbe d'abscisses $1$ et $3$.
\shortlines{1}
\item Recommencer avec les abscisses $1$ et $1{,}5$, puis $1$ et $1{,}1$.
\par\noindent\begin{tabularx}{\linewidth}{|>{\bfseries}p{40mm}|X|X|X|}
\hline
Second point d'abscisse & $3$ & $1{,}5$ & $1{,}1$ \\ \hline
Taux de variation & & & \\ \hline
\end{tabularx}
\item De quel nombre ces taux semblent-ils se rapprocher ? \rule{40mm}{0.32pt}
\item Écrire une phrase expliquant ce que ce nombre décrira en Première.
\worklines{1}
\end{enumerate}

\begin{successbox}
\textbf{Preuve attendue :} la relation de récurrence est écrite avec $u_{n+1}$ et $u_n$, et le taux de variation
est calculé comme un quotient de deux différences, dans le bon ordre.
\end{successbox}

\needspace{62mm}\par\vspace{3mm}
\phasehead{Phase 5 --- Avant l'évaluation}{70--75 min}{Cadrage méthodologique. Aucun contenu de l'évaluation n'est annoncé.}

\begin{methodbox}
\textbf{Cinq règles, et rien d'autre.} Ce cadrage ne porte sur aucun contenu de l'évaluation.
\begin{enumerate}
\item \textbf{Gérer le temps.} Trois parties : environ 10 min, 20 min et 15 min. Un regard sur l'horloge à la fin de chaque partie suffit.
\item \textbf{Justifier.} Écrire la relation, la propriété ou la règle avant le calcul. Un résultat seul ne montre pas ce que tu sais.
\item \textbf{Contrôler.} Avant de passer à la suite : ordre de grandeur, substitution, unité, ou vérification par une seconde voie.
\item \textbf{Lire jusqu'au bout.} Souligner la question réellement posée et les données effectivement fournies.
\item \textbf{Passer et revenir.} Une question laissée de côté puis reprise vaut mieux qu'une réponse écrite au hasard.
\end{enumerate}
\end{methodbox}

\begin{graybox}
\textbf{La certitude.} Elle est demandée à trois moments seulement, comme au test de positionnement de début de stage :
\conf\ . Elle n'entre pas dans la note. Elle sert à distinguer ce que tu sais de ce que tu crois savoir.
\end{graybox}

\subsectiontitle{Mon engagement d'avant-évaluation}
\par\vspace{1mm}
\noindent\textbf{Le contrôle que j'écrirai systématiquement :}\ \rule{\dimexpr\linewidth-76mm\relax}{0.32pt}
\par\vspace{3mm}
\noindent\textbf{La question que je traiterai en dernier :}\ \rule{\dimexpr\linewidth-68mm\relax}{0.32pt}
\par\vspace{1mm}

\begin{warningbox}\textbf{Fin du livret de travail.} L'évaluation finale est un document séparé, remis par l'enseignant à l'issue de ces 75 minutes.\end{warningbox}
\end{document}
